\documentclass[a4paper,12pt]{ctexart}
\usepackage{xcolor,graphicx,amsmath}
\usepackage[top=1.8cm, bottom=1.8cm, left=1.8cm, right=1.8cm]{geometry}
\usepackage{array,hyperref}
\hypersetup{colorlinks=true,linkcolor=blue!70!black}
\linespread{1.35}
\definecolor{defblue}{RGB}{0,80,160}\definecolor{thinkorange}{RGB}{230,120,20}
\definecolor{funboxbg}{RGB}{245,248,252}\definecolor{practicegreen}{RGB}{40,130,70}
\newenvironment{thinkbox}{\vspace{0.3cm}\begin{center}\begin{minipage}{0.88\textwidth}\colorbox{funboxbg}{\begin{minipage}{0.94\textwidth}\vspace{0.15cm}\textbf{\textcolor{thinkorange}{【 想一想 】}}}{\vspace{0.15cm}\end{minipage}}\end{minipage}\end{center}\vspace{0.3cm}}
\newenvironment{definition}[1]{\vspace{0.2cm}\begin{center}\begin{minipage}{0.88\textwidth}\fbox{\begin{minipage}{0.92\textwidth}\vspace{0.1cm}\textbf{\textcolor{defblue}{【 #1 】}}}{\vspace{0.1cm}\end{minipage}}\end{minipage}\end{center}\vspace{0.2cm}}
\newenvironment{example}{\vspace{0.2cm}\begin{center}\begin{minipage}{0.88\textwidth}\colorbox{examplebg}{\begin{minipage}{0.94\textwidth}\vspace{0.15cm}\textbf{\textcolor{orange!80!black}{【 例题 】}}}{\vspace{0.15cm}\end{minipage}}\end{minipage}\end{center}\vspace{0.2cm}}
\newenvironment{practice}{\vspace{0.3cm}\begin{center}\begin{minipage}{0.88\textwidth}\colorbox{funboxbg}{\begin{minipage}{0.94\textwidth}\vspace{0.15cm}\textbf{\textcolor{practicegreen}{【 练一练 】}}}{\vspace{0.15cm}\end{minipage}}\end{minipage}\end{center}\vspace{0.2cm}}
\newenvironment{figplaceholder}[1]{\vspace{0.2cm}\begin{center}\fbox{\begin{minipage}{0.82\textwidth}\centering\textcolor{blue!70!black}{\textbf{【插图位置】}}\\[0.2cm]#1\end{minipage}}\end{center}\vspace{0.2cm}}
\title{\textcolor{blue!70!black}{\Huge\textbf{生物4 · 生态与生物多样性}}}
\author{}\date{}
\begin{document}\maketitle\thispagestyle{empty}
\section*{\textcolor{blue!70!black}{致同学}}
这是生物系列的最后一本——我们把视野拉到"最大视角"：从个体之间的关系到整个生物圈的运转。生态学研究的是"谁吃谁、谁依赖谁、谁影响了谁"——它是一个复杂的网络，而人类是这个网络中影响力最大的节点。

\newpage
\section{\textcolor{orange!80!black}{第一课\quad 生物分类}}

\subsection*{为什么要分类？}
地球上有约870万种已经命名和近亿种未命名的生物——没有分类系统就是一团乱麻。瑞典科学家\textbf{林奈}创立了分类学。

\subsection*{七级分类系统}
\textbf{界门纲目科属种}——从大到小。种（物种）是最基本的分类单位。\\
\textbf{记忆口诀：}界门纲目科属种——"芥末刚摸蝌蚪肿"。

\subsection*{双名法——林奈的命名规则}
每个物种有一个\textbf{属名 + 种加词}（拉丁文，斜体）。如：\textit{Homo sapiens}（智人——人类）。属名大写，种加词小写。

\subsection*{五界系统（最常用）}
\textbf{原核生物界}（细菌/蓝藻）$\rightarrow$ \textbf{原生生物界}（草履虫/变形虫）$\rightarrow$ \textbf{真菌界}（酵母/蘑菇）$\rightarrow$ \textbf{植物界} $\rightarrow$ \textbf{动物界}。

\begin{thinkbox}
人类在分类上属于：动物界→脊索动物门→哺乳纲→灵长目→人科→人属→智人。你和你家的猫（动物界→脊索动物门→哺乳纲→食肉目→猫科→猫属→猫）——在"目"这个地方开始分岔。
\end{thinkbox}

\begin{practice}
1. 生物分类的七个等级从大到小是：\_\_\_\_\_。\\
2. 最基本（最小）的分类单位是\_\_\_\_\_。\\
3. 双名法中，\textit{Homo sapiens}中Homo是\_\_\_\_\_名，sapiens是\_\_\_\_\_。\\
4. 人类属于\_\_\_\_\_目。\\
5. （判断）"种"是最基本的分类单位。（\quad）\\
6. 五界系统中，细菌属于\_\_\_\_\_界。\\
7. 双名法的优点是可以避免\_\_\_\_\_混乱。\\
8. 在分类等级中，\textbf{科}比\textbf{目}的范畴\_\_\_\_\_（大/小）。
\end{practice}

\newpage
\section{\textcolor{orange!80!black}{第二课\quad 生物多样性}}

\subsection*{三个层次的多样性}
\textbf{物种多样性}：地球上不同种类的生物的丰富程度。热带雨林——面积只占全球陆地7\%——却包含了50\%以上的物种。\\
\textbf{遗传多样性（基因多样性）}：同一物种内不同个体之间的基因差异。人类个体之间DNA序列差异约0.1\%（约300万个碱基不同）。\\
\textbf{生态系统多样性}：不同的生境类型——森林/草原/湿地/海洋/沙漠。

\subsection*{中国是"巨大多样性"国家}
中国拥有哺乳动物超过500种（世界第一）、鸟类约1300种、植物约3万种。原因：国土纬度跨度大（北到寒温带、南到热带）、地形复杂（青藏高原/盆地/平原/丘陵）。是全球12个"巨大多样性国家"之一。

\begin{practice}
1. 生物多样性的三个层次：\_\_\_\_\_多样性、\_\_\_\_\_多样性和\_\_\_\_\_多样性。\\
2. （判断）生物多样性只指物种的丰富程度。（\quad）\\
3. 热带雨林面积仅占全球陆地7\%，却包含了世界\_\_\_\_\_\%以上的物种。\\
4. 中国生物多样性极其丰富，主要原因是\_\_\_\_\_跨度大和\_\_\_\_\_多样。\\
5. 遗传多样性指同一物种内个体之间的\_\_\_\_\_差异。\\
6. 人类个体之间DNA序列差异约\_\_\_\_\_。\\
7. （判断）生态系统多样性是生物多样性的最高层次。（\quad）\\
8. 为什么中国能成为生物多样性最丰富的国家之一？从纬度和地形两个角度分析。
\end{practice}

\newpage
\section{\textcolor{orange!80!black}{第三课\quad 种群的特征}}

\subsection*{种群——同一地方的同种个体}
\textbf{种群密度}：单位面积（或体积）内的个体数。如何调查？——用\textbf{样方法}（随机放置方格计数）或\textbf{标志重捕法}（捉一批标记后放回、再捉一批、按比例估算总数）。

\subsection*{种群的基本特征}
\textbf{出生率和死亡率}：决定了种群大小变化。出生率 > 死亡率 $\rightarrow$ 种群增长。\\
\textbf{迁入率和迁出率}：对城市人口影响极大（北京每年新增人口中很大比例来自迁入）。\\
\textbf{年龄结构}：年轻个体多——增长型；中老年多——衰退型；各年龄均衡——稳定型。

\subsection*{种群增长模型}
\textbf{J形曲线（理想条件）}：食物空间无限、没有天敌——种群指数增长。但自然界不可能无限。\\
\textbf{S形曲线（现实条件）}：种群增长到一定程度后增速减慢——达到环境容纳量\textbf{K值}（环境所能支持的最大种群数量）。K/2处种群增长率最大（这个点对合理捕捞/狩猎很重要——在K/2附近捕，种群恢复最快）。

\begin{example}
一个湖泊的最大鱼群容纳量为10000尾。问：鱼群数量在多少时捕捞效率最高？

\textbf{解：}K/2 = 5000尾——此时种群增长率最大。从5000尾捕到4000尾——种群能快速恢复到5000尾附近。若从9000尾捕到8000尾——增长慢，恢复时间长。所以合理捕捞应该控制在K/2附近。
\end{example}

\begin{practice}
1. 调查草地中蒲公英的数量宜用\_\_\_\_\_法。\\
2. 种群增长的S形曲线中，增长率最大时种群数量约为\_\_\_\_\_（填K或K/2）。\\
3. （判断）环境容纳量K值是不变的。（\quad）\\
4. 一座城市人口增长主要受\_\_\_\_\_和\_\_\_\_\_影响。\\
5. 年龄结构为增长型的种群，\_\_\_\_\_个体多。\\
6. （判断）标志重捕法适用于调查活动能力强的动物。（\quad）\\
7. 环境容纳量用字母\_\_\_\_\_表示。\\
8. J形增长曲线是在\_\_\_\_\_条件下出现的。
\end{practice}

\newpage
\section{\textcolor{orange!80!black}{第四课\quad 群落与种间关系}}

\subsection*{群落——同一区域的所有种群}
群落 = 所有生物的集合。不同的群落有不同的\textbf{物种组成}和\textbf{丰富度}。

\subsection*{种间关系}
\textbf{捕食}：一种生物吃另一种（狼吃羊）。捕食者和被捕食者互为影响——数量波动此消彼长。\\
\textbf{竞争}：两种生物争夺同一资源（水和食物）。竞争的结果——一个占优势，一个被排挤（竞争排斥原理）。\\
\textbf{寄生}：一种生物（寄生者）生活在另一种（宿主）体内或体表从宿主获取营养（蛔虫、跳蚤、菟丝子）。\\
\textbf{互利共生}：两者互相依赖、互相受益。地衣=真菌+藻类（真菌提供水和矿物质、藻类提供有机物）。根瘤菌和豆科植物（根瘤菌固氮→植物吸收；植物提供营养给根瘤菌）。

\begin{figplaceholder}
此处插入四种种间关系示意图（捕食/竞争/寄生/互利共生）（见 figure/requirements/bio4-ch1.txt 插图1）
\end{figplaceholder}

\subsection*{群落演替——"换班"的过程}
\textbf{初生演替}：在原本没有生物的地方（一块裸露的岩石）——最先地衣→苔藓→草本→灌木→森林——经过漫长的时间，变成稳定的群落。\\
\textbf{次生演替}：在原有植被被破坏但土壤还在的地方——更快恢复（如弃耕农田→杂草→灌木→森林）。

\begin{practice}
1. 两种生物为了同一资源而竞争——叫\_\_\_\_\_关系。\\
2. 地衣是真菌和藻类的\_\_\_\_\_关系。\\
3. （判断）一片荒漠中最初的生物定居过程叫次生演替。（\quad）\\
4. 捕食者和被捕食者数量波动的关系是\_\_\_\_\_。\\
5. （判断）两种生物的竞争结果总是两败俱伤。（\quad）\\
6. 寄生者从宿主获取\_\_\_\_\_。\\
7. 根瘤菌与豆科植物之间是\_\_\_\_\_关系。\\
8. 在原本没有生物的地方开始的演替叫\_\_\_\_\_演替。
\end{practice}

\newpage
\section{\textcolor{orange!80!black}{第五课\quad 生态系统的结构}}

\subsection*{什么是生态系统？}
生态系统 = 生物群落 + 无机环境。\textbf{范围可大可小}——一整个海洋是生态系统，一个池塘也是。全球最大的生态系统——\textbf{生物圈}。

\subsection*{四组成分}
\begin{enumerate}
  \item \textbf{非生物的物质和能量}：阳光、水、空气、温度、土壤。
  \item \textbf{生产者}：绿色植物（光合作用把无机物变有机物）——生态系统的"奠基者"。
  \item \textbf{消费者}：动物（直接或间接以生产者为食）。一级消费者（吃植物）、二级消费者（吃草食动物）、三级消费者（吃肉食动物）。
  \item \textbf{分解者}：细菌和真菌——把有机残体分解成无机物（矿物质）——回归土壤供生产者再利用。
\end{enumerate}

\subsection*{食物链与食物网}
\textbf{食物链}：草$\rightarrow$蝗虫$\rightarrow$青蛙$\rightarrow$蛇$\rightarrow$鹰——一条链。\\
\textbf{食物网}：多条食物链交织成网状——更接近真实的生态关系。食物网越复杂，生态系统越稳定（某一种生物数量变化不会引起整个系统崩溃）。

\begin{practice}
1. 生态系统的四组成分：\_\_\_\_\_、\_\_\_\_\_、\_\_\_\_\_、\_\_\_\_\_。\\
2. （判断）所有的动物都是消费者。（\quad）\\
3. 在食物网中，某种生物消失——食物网越\_\_\_\_\_（庞大/简单）越容易崩溃。\\
4. 生产者主要是指\_\_\_\_\_，它们能将光能转化为\_\_\_\_\_能。\\
5. 分解者主要是指\_\_\_\_\_和\_\_\_\_\_。\\
6. （判断）食物网越复杂，生态系统越稳定。（\quad）\\
7. 非生物成分包括阳光、\_\_\_\_\_、\_\_\_\_\_等。\\
8. 三级消费者以\_\_\_\_\_（初级/次级）消费者为食。
\end{practice}

\newpage
\section{\textcolor{orange!80!black}{第六课\quad 生态系统的功能}}

\subsection*{能量流动}
能量从生产者（植物）流入消费者——沿食物链单向传递。每个营养级的能量大部分散失（以体热等形式）——大约只有\textbf{10\%-20\%}传递到下一营养级。这就解释了为什么食物链一般不超过4-5级——能量"不够用了"。

\textbf{生态金字塔}：植物最多（第一营养级）$\rightarrow$ 草食动物 $\rightarrow$ 肉食动物——越往上生物量越少。这就是自然界的"收入分配"法则。

\subsection*{物质循环——碳循环}
碳在无机环境和生物群落之间循环：\\
CO$_2$（大气）$\xrightarrow{\text{光合作用}}$ 有机物（植物体）$\xrightarrow{\text{呼吸/分解}}$ CO$_2$（回到大气）。\\
化石燃料（煤/石油/天然气）的燃烧——把亿万年储存的碳快速释放到大气中——加剧\textbf{温室效应}。

\subsection*{生物富集}
有毒物质（如DDT、汞）在食物链中传递时，在最高营养级体内\textbf{浓度越来越高}——这叫\textbf{生物富集}。鹰体内的DDT浓度可能是水中浓度的数百万倍。所以一旦环境被重金属污染——站在食物链顶端的人类（和肉食动物）受害最深。

\begin{example}
一亩水稻田年产稻谷500kg——吃掉这些稻谷的人最多可获得多少kg的生物量？（假设能量传递效率10\%）

\textbf{解：}水稻→人：一级传递。人获得的生物量$\approx$500kg$\times$10\%=50kg。实际上人吃的不仅是稻谷——还有家畜等——但总体上能量越传递越少。
\end{example}

\begin{practice}
1. 能量沿食物链传递的效率大约是\_\_\_\_\%。\\
2. （判断）碳循环中，大气中的CO$_2$通过光合作用进入生物群落。（\quad）\\
3. 生物富集使最高营养级体内的有毒物质浓度\_\_\_\_\_（最高/最低）。\\
4. 温室效应加剧的主要原因是\_\_\_\_\_。\\
5. 能量沿食物链传递的方向是\_\_\_\_\_（单向/循环）的。\\
6. 生态金字塔中数量最多的营养级通常是\_\_\_\_\_。\\
7. （判断）碳在无机环境和生物群落间循环往复。（\quad）\\
8. 生物富集使\_\_\_\_\_（低/高）营养级体内毒物浓度最高。
\end{practice}

\newpage
\section{\textcolor{orange!80!black}{第七课\quad 生态系统的稳定性}}

\subsection*{抵抗力 vs 恢复力}
\textbf{抵抗力稳定性}：抵抗外界干扰保持平衡的能力。生态系统越复杂（食物网越密），抵抗力越强——热带雨林。\\
\textbf{恢复力稳定性}：被破坏后恢复到平衡的能力。简单生态系统容易恢复——荒漠、苔原。\\
两者通常\textbf{负相关}——雨林一旦破坏极难恢复（恢复力低）；荒漠恢复力高（但抵抗力低）。

\subsection*{生态系统受到的影响}
\textbf{自然干扰}：火灾、洪水、虫害——大多数生态系统可以自然恢复。\\
\textbf{人为干扰}：森林砍伐、过度捕捞、环境污染——往往超过生态系统的调节能力（因为干扰持续且极强）。

\begin{practice}
1. 生态系统的稳定性分为\_\_\_\_\_稳定性和\_\_\_\_\_稳定性。\\
2. （判断）热带雨林的恢复力稳定性很高。（\quad）\\
3. 极地苔原生态系统很脆弱——主要是因为\_\_\_\_\_稳定性低。\\
4. （判断）生态系统的抵抗力稳定性和恢复力稳定性通常正相关。（\quad）\\
5. 热带雨林的抵抗力稳定性\_\_\_\_\_（高/低）。\\
6. 荒漠生态系统的恢复力稳定性\_\_\_\_\_（高/低）。\\
7. 生态系统的自我调节能力有一定\_\_\_\_\_。
\end{practice}

\newpage
\section{\textcolor{orange!80!black}{第八课\quad 人类活动与环境问题}}

\subsection*{主要环境问题}
\textbf{全球变暖}：CO$_2$等温室气体增多——地球平均温度上升——冰川融化——海平面上升——极端天气频发——对北京的影响：夏季更热、暴雨更极端。\\
\textbf{臭氧层破坏}：氟利昂等化学物质破坏臭氧层——紫外线增强——皮肤癌和白内障风险增加。\\
\textbf{酸雨}：SO$_2$和氮氧化物排放$\rightarrow$形成硫酸和硝酸——使土壤酸化、腐蚀建筑。\\
\textbf{荒漠化}：过度放牧+过度开垦→沙漠面积扩大。北京春天的沙尘暴就是荒漠化的后果之一。

\subsection*{可持续发展的思路}
\textbf{绿水青山就是金山银山。}经济发展和环境保护不是"二选一"——而是长期来看必须兼顾。具体的思路："双碳"目标（碳达峰、碳中和）、退耕还林/还草、垃圾分类与循环利用。

\begin{practice}
1. 全球变暖的主要原因是\_\_\_\_\_气体增多。\\
2. 臭氧层破坏会导致\_\_\_\_\_辐射增强。\\
3. （判断）酸雨的主要成分是硫酸和硝酸。（\quad）\\
4. 北京春天的沙尘暴与\_\_\_\_\_化有关。\\
5. 臭氧层被破坏的主要原因是\_\_\_\_\_的排放。\\
6. （判断）酸雨是自然现象，与人类活动无关。（\quad）\\
7. "双碳"目标是指\_\_\_\_\_和\_\_\_\_\_。\\
8. 垃圾分类属于\_\_\_\_\_（资源循环/垃圾处理）理念。
\end{practice}

\newpage
\section{\textcolor{orange!80!black}{第九课\quad 生态保护实践}}

\subsection*{自然保护区}
设立自然保护区——保护珍稀物种及其栖息地。北京周边的自然保护区：百花山、松山、喇叭沟门。

\subsection*{生物多样性的保护}
\textbf{就地保护}：在原有生境中保护（建立自然保护区）。\\
\textbf{迁地保护}：在人工条件下保护（植物园、动物园、种子库）。\\
\textbf{法律法规}：禁止非法猎捕、贸易濒危物种。

\subsection*{每个人都可以做的事}
\begin{itemize}
  \item 少用一次性塑料制品（塑料分解需要几百年）
  \item 节约用水用电（背后是能源消耗）
  \item 垃圾分类：北京已经全面实行
  \item 多步行少开车
  \item 不购买野生动物制品
\end{itemize}

\begin{thinkbox}
你在自然故事里学过的所有"为什么"——天空为什么蓝、水循环怎么走、沙尘暴从哪来、恐龙为什么灭绝——你现在有了生物学、物理、化学、地理的知识框架来重新理解它们。这就是"通识教育"的终点和起点。
\end{thinkbox}

\begin{practice}
1. 保护生物多样性最有效的措施是\_\_\_\_\_保护。\\
2. （判断）北京已经全面实行垃圾分类。（\quad）\\
3. 自然保护区属于\_\_\_\_\_（就地/迁地）保护。\\
4. 建立种子库属于\_\_\_\_\_保护。\\
5. （判断）自然保护区只能保护动物，不能保护植物。（\quad）\\
6. 北京百花山自然保护区属于\_\_\_\_\_（就地/迁地）保护。\\
7. 保护生物多样性最根本的措施是保护\_\_\_\_\_。\\
8. 不购买野生动物制品属于\_\_\_\_\_（个人行动/法律手段）。
\end{practice}

\vspace{0.5cm}
\begin{center}\rule{0.7\textwidth}{1pt}\vspace{0.4cm}
{\large\textcolor{blue!70!black}{\textbf{—— 生物系列（4本）完 ——}}}
\end{center}

\newpage
\section*{\textcolor{defblue}{参考答案}}

\subsection*{第一课}
1. 界、门、纲、目、科、属、种\\
2. 种（物种）\\
3. 属、种加词\\
4. 灵长\\
5. （√）\\
6. 原核生物\\
7. 命名\\
8. 小

\subsection*{第二课}
1. 物种、遗传（基因）、生态系统\\
2. （×）包括三个层次\\
3. 50\\
4. 纬度、地形\\
5. 基因（DNA序列）\\
6. 0.1\%（约300万个碱基）\\
7. （√）\\
8. 中国纬度跨度大（北到寒温带、南到热带），地形复杂（青藏高原/盆地/平原/丘陵），多样的气候和地貌为不同生物提供了丰富的栖息地。

\subsection*{第三课}
1. 样方\\
2. K/2\\
3. （×）K值随环境条件改变\\
4. 迁入率、迁出率\\
5. 年轻（幼年）\\
6. （√）\\
7. K\\
8. 理想

\subsection*{第四课}
1. 竞争\\
2. 互利共生\\
3. （×）属于初生演替\\
4. 此消彼长（相互制约）\\
5. （×）竞争一方占优势，另一方被排挤\\
6. 营养\\
7. 互利共生\\
8. 初生

\subsection*{第五课}
1. 非生物的物质和能量、生产者、消费者、分解者\\
2. （×）秃鹫、蜣螂等属于分解者\\
3. 简单\\
4. 绿色植物（生产者）、化学\\
5. 细菌、真菌\\
6. （√）\\
7. 水、空气（温度、土壤等）\\
8. 次级

\subsection*{第六课}
1. 10\%--20\%\\
2. （√）\\
3. 最高\\
4. 化石燃料大量燃烧（CO$_2$排放）\\
5. 单向\\
6. 生产者（第一营养级）\\
7. （√）\\
8. 高

\subsection*{第七课}
1. 抵抗力、恢复力\\
2. （×）雨林恢复力稳定性低\\
3. 抵抗力\\
4. （×）两者通常负相关\\
5. 高\\
6. 高\\
7. 限度

\subsection*{第八课}
1. 温室（CO$_2$等）\\
2. 紫外线（UV）\\
3. （√）\\
4. 荒漠\\
5. 氟利昂（氟氯烃）\\
6. （×）人类活动排放SO$_2$和氮氧化物是主因\\
7. 碳达峰、碳中和\\
8. 资源循环

\subsection*{第九课}
1. 就地\\
2. （√）\\
3. 就地\\
4. 迁地\\
5. （×）自然保护区保护整个生态系统\\
6. 就地\\
7. 栖息地（生态环境）\\
8. 个人行动

\end{document}
